% Practice list B --- topics of Listas 3-4 (Kurlat chs. 6-7), GRADUATE LEVEL (style guide §8).
% \solucoesfalse -> blank list ; \solucoestrue -> answer key + rubric.
% Compile: pdflatex (2x); check: pdffonts (all Type 1). Every number: practice-L3-L4-check.py
\documentclass[11pt,a4paper]{article}
% Shared preamble of the practice lists (practice-L*.tex), adapted from lista-exam.tex.
% Each list sets \solucoesfalse / \solucoestrue and \listatema right after % Shared preamble of the practice lists (practice-L*.tex), adapted from lista-exam.tex.
% Each list sets \solucoesfalse / \solucoestrue and \listatema right after % Shared preamble of the practice lists (practice-L*.tex), adapted from lista-exam.tex.
% Each list sets \solucoesfalse / \solucoestrue and \listatema right after \input{practice-preamble}.
\usepackage[utf8]{inputenc}
\usepackage[T1]{fontenc}
\usepackage{lmodern}   % Type 1 vector fonts (replaces PK bitmaps)
\usepackage{cmap}      % ToUnicode CMap -> copyable text and accents
\usepackage[english]{babel}
\usepackage{amsmath,amssymb,amsfonts}
\usepackage[a4paper,margin=2.2cm]{geometry}
\usepackage{enumitem}
\usepackage{fancyhdr}
\usepackage{comment}
\usepackage{xcolor}
\usepackage{microtype}
\usepackage{booktabs}
\usepackage{tikz}

\newif\ifsolucoes
\newcommand{\disciplina}{Macroeconomics I --- MPE Insper}
\newcommand{\espacoresol}{4.5cm}   % height of the answer space in the blank list
% The sol environment must be decided after the list sets its toggle.
\AtBeginDocument{\ifsolucoes\includecomment{sol}\else\excludecomment{sol}\fi}

\pagestyle{fancy}\fancyhf{}
\lhead{\footnotesize\disciplina}
\rhead{\footnotesize Practice list\ifsolucoes{} --- Solutions and rubric\fi}
\cfoot{\footnotesize\thepage}
\renewcommand{\headrulewidth}{0.3pt}
\setlength{\headheight}{14pt}
\setlist{topsep=2pt,itemsep=2pt}
\newcounter{questao}

\newcommand{\bloco}[3]{\bigskip\noindent
  {\large\bfseries\color{blue!55!black}Block #1 --- #2}\par\noindent\hfill\textit{(#3)}\par
  \noindent\rule{\linewidth}{0.3pt}\smallskip}
\newcommand{\questao}[1]{\refstepcounter{questao}\medskip\noindent
  \textbf{Question \thequestao.}~#1\par}
\newenvironment{alternativas}{\begin{enumerate}[label=\Alph*.,leftmargin=2.2em,topsep=1pt]}{\end{enumerate}}
\newcommand{\gabarito}[1]{\ifsolucoes\par\noindent\textbf{Answer:}~#1\fi}
\newcommand{\resolespaco}{\ifsolucoes\else\par\textit{Solution:}\vspace{\espacoresol}\par\fi}
\newcommand{\passo}[2]{\par\smallskip\noindent\textbf{Step --- #1.}~#2}
\newcommand{\intuicao}[1]{\par\smallskip\noindent\textit{Economic intuition.}~#1}
\newcommand{\refbib}[1]{\par\smallskip\noindent{\footnotesize\color{black!60}\textbf{Ref:} #1}}
\newcommand{\rubrica}[1]{\par\smallskip\noindent\textbf{Grading rubric.}\begin{itemize}[leftmargin=1.4em,topsep=1pt]#1\end{itemize}}

\newcommand{\cabecalho}{%
\begin{center}
  {\Large\bfseries \disciplina}\\[2pt]
  {\large Practice list: \listatema}\ifsolucoes\\[2pt]\textit{Solutions and grading rubric}\fi
\end{center}
\vspace{2pt}
{\footnotesize\noindent\textbf{Instructions.} Multiple choice: mark the single correct option (no
penalty for wrong answers). \emph{True/False}: every 2 wrong items cancel 1 right one
(score $=\max\{0,\text{right}-\text{wrong}/2\}\times$ item value); blank items do not count.
Discursive questions: show every step. Numbers are illustrative and chosen to give clean answers.\par}}
.
\usepackage[utf8]{inputenc}
\usepackage[T1]{fontenc}
\usepackage{lmodern}   % Type 1 vector fonts (replaces PK bitmaps)
\usepackage{cmap}      % ToUnicode CMap -> copyable text and accents
\usepackage[english]{babel}
\usepackage{amsmath,amssymb,amsfonts}
\usepackage[a4paper,margin=2.2cm]{geometry}
\usepackage{enumitem}
\usepackage{fancyhdr}
\usepackage{comment}
\usepackage{xcolor}
\usepackage{microtype}
\usepackage{booktabs}
\usepackage{tikz}

\newif\ifsolucoes
\newcommand{\disciplina}{Macroeconomics I --- MPE Insper}
\newcommand{\espacoresol}{4.5cm}   % height of the answer space in the blank list
% The sol environment must be decided after the list sets its toggle.
\AtBeginDocument{\ifsolucoes\includecomment{sol}\else\excludecomment{sol}\fi}

\pagestyle{fancy}\fancyhf{}
\lhead{\footnotesize\disciplina}
\rhead{\footnotesize Practice list\ifsolucoes{} --- Solutions and rubric\fi}
\cfoot{\footnotesize\thepage}
\renewcommand{\headrulewidth}{0.3pt}
\setlength{\headheight}{14pt}
\setlist{topsep=2pt,itemsep=2pt}
\newcounter{questao}

\newcommand{\bloco}[3]{\bigskip\noindent
  {\large\bfseries\color{blue!55!black}Block #1 --- #2}\par\noindent\hfill\textit{(#3)}\par
  \noindent\rule{\linewidth}{0.3pt}\smallskip}
\newcommand{\questao}[1]{\refstepcounter{questao}\medskip\noindent
  \textbf{Question \thequestao.}~#1\par}
\newenvironment{alternativas}{\begin{enumerate}[label=\Alph*.,leftmargin=2.2em,topsep=1pt]}{\end{enumerate}}
\newcommand{\gabarito}[1]{\ifsolucoes\par\noindent\textbf{Answer:}~#1\fi}
\newcommand{\resolespaco}{\ifsolucoes\else\par\textit{Solution:}\vspace{\espacoresol}\par\fi}
\newcommand{\passo}[2]{\par\smallskip\noindent\textbf{Step --- #1.}~#2}
\newcommand{\intuicao}[1]{\par\smallskip\noindent\textit{Economic intuition.}~#1}
\newcommand{\refbib}[1]{\par\smallskip\noindent{\footnotesize\color{black!60}\textbf{Ref:} #1}}
\newcommand{\rubrica}[1]{\par\smallskip\noindent\textbf{Grading rubric.}\begin{itemize}[leftmargin=1.4em,topsep=1pt]#1\end{itemize}}

\newcommand{\cabecalho}{%
\begin{center}
  {\Large\bfseries \disciplina}\\[2pt]
  {\large Practice list: \listatema}\ifsolucoes\\[2pt]\textit{Solutions and grading rubric}\fi
\end{center}
\vspace{2pt}
{\footnotesize\noindent\textbf{Instructions.} Multiple choice: mark the single correct option (no
penalty for wrong answers). \emph{True/False}: every 2 wrong items cancel 1 right one
(score $=\max\{0,\text{right}-\text{wrong}/2\}\times$ item value); blank items do not count.
Discursive questions: show every step. Numbers are illustrative and chosen to give clean answers.\par}}
.
\usepackage[utf8]{inputenc}
\usepackage[T1]{fontenc}
\usepackage{lmodern}   % Type 1 vector fonts (replaces PK bitmaps)
\usepackage{cmap}      % ToUnicode CMap -> copyable text and accents
\usepackage[english]{babel}
\usepackage{amsmath,amssymb,amsfonts}
\usepackage[a4paper,margin=2.2cm]{geometry}
\usepackage{enumitem}
\usepackage{fancyhdr}
\usepackage{comment}
\usepackage{xcolor}
\usepackage{microtype}
\usepackage{booktabs}
\usepackage{tikz}

\newif\ifsolucoes
\newcommand{\disciplina}{Macroeconomics I --- MPE Insper}
\newcommand{\espacoresol}{4.5cm}   % height of the answer space in the blank list
% The sol environment must be decided after the list sets its toggle.
\AtBeginDocument{\ifsolucoes\includecomment{sol}\else\excludecomment{sol}\fi}

\pagestyle{fancy}\fancyhf{}
\lhead{\footnotesize\disciplina}
\rhead{\footnotesize Practice list\ifsolucoes{} --- Solutions and rubric\fi}
\cfoot{\footnotesize\thepage}
\renewcommand{\headrulewidth}{0.3pt}
\setlength{\headheight}{14pt}
\setlist{topsep=2pt,itemsep=2pt}
\newcounter{questao}

\newcommand{\bloco}[3]{\bigskip\noindent
  {\large\bfseries\color{blue!55!black}Block #1 --- #2}\par\noindent\hfill\textit{(#3)}\par
  \noindent\rule{\linewidth}{0.3pt}\smallskip}
\newcommand{\questao}[1]{\refstepcounter{questao}\medskip\noindent
  \textbf{Question \thequestao.}~#1\par}
\newenvironment{alternativas}{\begin{enumerate}[label=\Alph*.,leftmargin=2.2em,topsep=1pt]}{\end{enumerate}}
\newcommand{\gabarito}[1]{\ifsolucoes\par\noindent\textbf{Answer:}~#1\fi}
\newcommand{\resolespaco}{\ifsolucoes\else\par\textit{Solution:}\vspace{\espacoresol}\par\fi}
\newcommand{\passo}[2]{\par\smallskip\noindent\textbf{Step --- #1.}~#2}
\newcommand{\intuicao}[1]{\par\smallskip\noindent\textit{Economic intuition.}~#1}
\newcommand{\refbib}[1]{\par\smallskip\noindent{\footnotesize\color{black!60}\textbf{Ref:} #1}}
\newcommand{\rubrica}[1]{\par\smallskip\noindent\textbf{Grading rubric.}\begin{itemize}[leftmargin=1.4em,topsep=1pt]#1\end{itemize}}

\newcommand{\cabecalho}{%
\begin{center}
  {\Large\bfseries \disciplina}\\[2pt]
  {\large Practice list: \listatema}\ifsolucoes\\[2pt]\textit{Solutions and grading rubric}\fi
\end{center}
\vspace{2pt}
{\footnotesize\noindent\textbf{Instructions.} Multiple choice: mark the single correct option (no
penalty for wrong answers). \emph{True/False}: every 2 wrong items cancel 1 right one
(score $=\max\{0,\text{right}-\text{wrong}/2\}\times$ item value); blank items do not count.
Discursive questions: show every step. Numbers are illustrative and chosen to give clean answers.\par}}

\solucoesfalse                 % <<<< \solucoestrue for the ANSWER KEY
\newcommand{\listatema}{Consumption, Saving and Labor Markets --- graduate level (Listas 3--4)}
\renewcommand{\espacoresol}{6cm}

\begin{document}
\cabecalho

% =============================================================================
\bloco{1}{Consumption and saving}{Kurlat ch.~6; Romer ch.~8 \,\textbullet\, 25 points}

\questao{(10 points) Objective items.}

\textbf{a)} (3 points) A household lives two periods with $u(c_1)+\beta u(c_2)$, $u=c^{1-\sigma}/(1-\sigma)$,
$\sigma=2$, $\beta=0.95$, and income $(y_1,y_2)$. It earns $r_s=2\%$ (real) on deposits but pays
$r_b=30\%$ (real) on consumer credit. For which values of $y_2/y_1$ does it neither borrow nor lend?
\begin{alternativas}
\item $[0.984,\ 1.111]$
\item $[0.969,\ 1.235]$
\item $[1.020,\ 1.300]$
\item $[0.950,\ 1.000]$
\end{alternativas}
\gabarito{A}
\begin{sol}
\passo{kinked budget}{The budget line has slope $-(1+r_s)$ to the right of the endowment (saving) and
$-(1+r_b)$ to its left (borrowing). At the kink $c_t=y_t$.}
\passo{optimality at the kink}{The household stays at the kink iff its MRS lies between the two gross
rates: $1+r_s\le\dfrac{u'(y_1)}{\beta u'(y_2)}=\dfrac{(y_2/y_1)^{\sigma}}{\beta}\le1+r_b$, i.e.\
$y_2/y_1\in[(\beta(1+r_s))^{1/\sigma},(\beta(1+r_b))^{1/\sigma}]=[0.984,1.111]$. Option B is the
log-utility ($\sigma=1$) answer.}
\refbib{Kurlat (2020), \S6.2, pp.~105--117, and Exercise 6.6 (kinked budget sets).}
\end{sol}

\medskip
\textbf{b)} (3 points) Three periods, $r=0$, wealth $W=100$, log utility. The self at date $t$ maximizes
$\ln c_t+\beta\sum_{s>t}\ln c_s$ with $\beta=1/2$ (present bias, no long-run discounting). Self 0
anticipates how self 1 will behave. What share of consumption \emph{in every period} would self 0
give up to commit future selves to its own preferred plan?
\begin{alternativas}
\item $0\%$
\item $2.9\%$
\item $5.9\%$
\item $1.5\%$
\end{alternativas}
\gabarito{B}
\begin{sol}
\passo{self 1}{$\max\ln c_1+\beta\ln c_2$ s.t.\ $c_1+c_2=W_1$: $c_1=W_1/(1+\beta)$, $c_2=\beta W_1/(1+\beta)$.}
\passo{self 0}{$\max\ln c_0+\beta[\ln c_1+\ln c_2]=\ln c_0+2\beta\ln W_1+\text{const}$: $c_0=W/(1+2\beta)=50$,
so $W_1=50$, $c_1=33.3$, $c_2=16.7$. Self 0 would prefer $c_1=c_2=25$.}
\passo{value of commitment}{Gain $\beta[2\ln25-\ln33.3-\ln16.7]=0.0589$ utils. Cutting all three
consumptions by a share $x$ costs $(1+2\beta)\ln\frac1{1-x}$, so $x=1-e^{-0.0589/2}=2.9\%$.
Option A confuses ``naive = sophisticated under log'' (item c) with ``no gain from commitment''.}
\refbib{Romer (2012), \S8.6 and problem 8.15, pp.~397--404; Kurlat (2020), \S6.4, pp.~120--121.}
\end{sol}

\medskip
\textbf{c)} (2 points, T or F) In (b), a sophisticated self 0 consumes less at date 0 than a naive self 0,
who believes future selves will follow its plan, because it anticipates their overconsumption.
\gabarito{F}
\begin{sol}
\passo{naive}{Believing $c_1=c_2=W_1/2$, the naive self solves $\max\ln c_0+2\beta\ln(W_1/2)$: $c_0=W/(1+2\beta)$, the same.}
\passo{why}{Under log utility, future consumption is proportional to $W_1$ whatever the rule, so the
misbehavior changes continuation utility by a constant that does not depend on $W_1$ and leaves the
margin untouched. With $\sigma\ne1$ the two differ (the check script verifies that the sophisticate
consumes less at $\sigma=2$ and more at $\sigma=1/2$); log is the knife edge.}
\refbib{Romer (2012), problem 8.15.}
\end{sol}

\medskip
\textbf{d)} (2 points, T or F) A household with $y_1=0$ and $y_2>0$ that can borrow at $r$ lowers $c_1$ when
$r$ rises, whatever $\sigma>0$.
\gabarito{T}
\begin{sol}
\passo{closed form}{From the Euler equation $c_2=[\beta(1+r)]^{1/\sigma}c_1$ and the budget
$c_1+c_2/(1+r)=y_2/(1+r)$: $c_1=\dfrac{y_2}{(1+r)+\beta^{1/\sigma}(1+r)^{1/\sigma}}$; the denominator rises with
$r$ for every $\sigma>0$. Substitution, income (a borrower is poorer) and human-wealth effects all
lower $c_1$.}
\refbib{Kurlat (2020), \S6.2; Romer (2012), \S8.4, pp.~380--384.}
\end{sol}

% -----------------------------------------------------------------------------
\questao{(15 points) Discursive. Permanent income over an infinite horizon. A household maximizes
$\sum_{t\ge0}\beta^tu(c_t)$, $u=c^{1-\sigma}/(1-\sigma)$, subject to $c_t+a_{t+1}=y_t+(1+r)a_t$, $a_0$ given,
and the no-Ponzi condition $\lim_{T\to\infty}a_{T+1}/(1+r)^T\ge0$. Let $R\equiv1+r$ and $H_t\equiv\sum_{j\ge0}y_{t+j}R^{-j}$.}

\textbf{a)} (3 points) Derive the intertemporal budget constraint and the consumption function
$c_0=(1-\gamma)(Ra_0+H_0)$, $\gamma\equiv\beta^{1/\sigma}R^{1/\sigma-1}$. State the condition for a solution and show that
with log utility the propensity to consume out of wealth does not depend on $r$.
\resolespaco
\begin{sol}
\passo{budget}{Multiply the flow constraint by $R^{-t}$: $R^{-t}c_t+R^{-t}a_{t+1}=R^{-t}y_t+R^{-(t-1)}a_t$. Summing
from $0$ to $T$ telescopes the assets:
\[\sum_{t=0}^TR^{-t}c_t+R^{-T}a_{T+1}=\sum_{t=0}^TR^{-t}y_t+Ra_0.\]
At an optimum the no-Ponzi condition binds ($\lim R^{-T}a_{T+1}=0$; a positive limit would be
unspent wealth). Letting $T\to\infty$: $\sum_tR^{-t}c_t=Ra_0+H_0$.}
\passo{Euler}{Lagrangian with multiplier $\Lambda$ on the budget: $\beta^tu'(c_t)=\Lambda R^{-t}$, so
$u'(c_t)=\beta Ru'(c_{t+1})$ and $c_t=c_0(\beta R)^{t/\sigma}$.}
\passo{consumption function}{Substituting: $c_0\sum_t(\beta R)^{t/\sigma}R^{-t}=c_0\sum_t\gamma^t=c_0/(1-\gamma)$, which needs
$\gamma<1$ (otherwise lifetime utility is unbounded and no optimum exists). Hence
$c_0=(1-\gamma)(Ra_0+H_0)$.}
\passo{log case}{$\sigma=1$: $\gamma=\beta$, so the propensity to consume out of total wealth is $1-\beta$,
independent of $r$ (income and substitution effects of $r$ cancel; $r$ still changes $H_0$).}
\refbib{Kurlat (2020), \S6.3, eqs.\ (6.3.6)--(6.3.9), pp.~117--120; Romer (2012), \S8.1.}
\rubrica{
\item Telescoping to the intertemporal constraint, with the role of no-Ponzi: 1 pt.
\item Euler equation: 1 pt.
\item Consumption function and $\gamma<1$: 0.5 pt.
\item Log case: 0.5 pt.
}
\end{sol}

\textbf{b)} (4 points) Let $\beta R=1$. Show that $c_t=\frac rR(Ra_t+H_t)$ and prove Campbell's
``saving for a rainy day'' equation, $s_t\equiv ra_t+y_t-c_t=-\sum_{j\ge1}R^{-j}\Delta y_{t+j}$.
\resolespaco
\begin{sol}
\passo{flat consumption}{$\beta R=1$ gives $\gamma=R^{-1}$, $1-\gamma=r/R$, and $c$ constant; applying (a) from
date $t$: $c_t=\frac rR(Ra_t+H_t)=ra_t+\frac rRH_t$.}
\passo{saving}{$s_t=ra_t+y_t-c_t=y_t-\frac rRH_t$.}
\passo{summation by parts}{Expand the right-hand side of Campbell's equation:
\[-\sum_{j\ge1}R^{-j}(y_{t+j}-y_{t+j-1})=-\sum_{j\ge1}R^{-j}y_{t+j}+R^{-1}\sum_{k\ge0}R^{-k}y_{t+k}
=-(H_t-y_t)+\frac{H_t}{R}=y_t-\frac rRH_t,\]
which is $s_t$ (assuming $R^{-j}y_{t+j}\to0$).}
\intuicao{Households save when they expect income to fall and dissave when they expect it to rise:
saving is a forecast of income declines, and consumption reacts to news, not to current income.}
\refbib{Romer (2012), \S8.2--8.3 (deterministic version).}
\rubrica{
\item Flat consumption and the annuity formula: 1 pt.
\item $s_t=y_t-\frac rRH_t$: 1 pt.
\item Summation by parts, every line shown: 2 pts. Verifying on an example only: max.\ 1 pt.
}
\end{sol}

\textbf{c)} (4 points) Brazil's 2019 pension reform, stylized. $r=4\%$, $\beta R=1$, $a_0=0$; earnings are
$100$ for $t=0,\dots,29$ and the pension is $60$ forever from $t=30$. At $t=0$ a reform is announced that
cuts the pension to $45$. Compute $\Delta c_0$ with the consumption function and $\Delta s_0$ with
Campbell's equation. Compare with the response to a one-off transfer and to a permanent raise of
the same size, and give the limit of $\Delta c_0$ as $r\to0$ (keeping $\beta R=1$).
\resolespaco
\begin{sol}
\passo{wealth effect}{$\Delta H_0=-15\sum_{t\ge30}R^{-t}=-15R^{-30}\frac{R}{r}$, so
$\Delta c_0=\frac rR\Delta H_0=-15R^{-30}=-15\cdot0.3083=-4.62$.}
\passo{Campbell}{Only $\Delta y_{30}$ changes, by $-15$: $\Delta s_0=-R^{-30}(-15)=+4.62=-\Delta c_0$. Consistent.}
\passo{benchmarks}{A one-off transfer $X$ at $t=0$: $\Delta c_0=\frac rRX=0.038X$. A permanent raise $X$:
$\Delta H_0=XR/r$, $\Delta c_0=X$.}
\passo{limit}{$-15R^{-30}\to-15$ as $r\to0$: a fully patient household cuts consumption today by the
whole future cut, because the infinite future weighs as much as the present.}
\intuicao{News about income 30 years ahead moves consumption today by 4.6\% of current earnings;
the pension reform raised private saving before any pension was paid.}
\refbib{Kurlat (2020), \S6.3; Romer (2012), \S8.2.}
\rubrica{
\item $\Delta H_0$ and $\Delta c_0$: 1.5 pts.
\item Campbell check: 1 pt.
\item Benchmarks: 1 pt. Treating the cut as transitory: max.\ 2 pts in the item.
\item Limit: 0.5 pt.
}
\end{sol}

\textbf{d)} (4 points) Now $a_{t+1}\ge0$ for all $t$ (no borrowing), log utility, $r=4\%$, $\beta R=1$, $a_0=0$;
earnings are $50$ for $t\le9$ and $100$ from $t=10$. Show the unconstrained plan is infeasible. Prove,
with the Kuhn--Tucker conditions, that $c_t=y_t$ is optimal and find which multipliers are positive.
Compute the welfare cost of the constraint (the share of the unconstrained consumption path the
household would accept to give up) and the propensity to consume a small transfer received at
$t=0$ and at $t=9$.
\resolespaco
\begin{sol}
\passo{infeasible}{Unconstrained $c=\frac rRH_0=50+50R^{-10}=83.78>50=y_0$, so $a_1<0$.}
\passo{Lagrangian}{$\mathcal L=\sum_t\beta^t\ln c_t+\sum_t\Lambda_t[y_t+Ra_t-c_t-a_{t+1}]+\sum_t\mu_{t+1}a_{t+1}$,
$\mu\ge0$, $\mu_{t+1}a_{t+1}=0$. FOCs: $\beta^t/c_t=\Lambda_t$ and $-\Lambda_t+R\Lambda_{t+1}+\mu_{t+1}=0$, so
$\mu_{t+1}=\beta^t\big(\frac1{c_t}-\frac1{c_{t+1}}\big)$ using $\beta R=1$.}
\passo{verification}{At $c_t=y_t$, $a_t=0$: $\mu_{t+1}=0$ for $t\le8$ and $t\ge10$ (consumption flat);
$\mu_{10}=\beta^9(\frac1{50}-\frac1{100})>0$. All multipliers are non-negative and complementary
slackness holds; the objective is concave and the constraints linear, so the KKT point is optimal.}
\passo{welfare}{$U^c=\sum_{t<10}\beta^t\ln50+\frac{\beta^{10}}{1-\beta}\ln100$, $U^u=\frac{\ln83.78}{1-\beta}$. The cost
$x$ solves $\frac{\ln(1-x)}{1-\beta}=U^c-U^u$: $x=4.68\%$ of consumption forever.}
\passo{propensities}{At $t=9$ the household wants to borrow against $t=10$ ($\mu_{10}>0$): a small
transfer is consumed at once, MPC $=1$. At $t=0$ the Euler equation holds with equality between
dates $0,\dots,9$, so the transfer is spread over those ten dates:
$\Delta c_0=X\frac{r/R}{1-R^{-10}}=0.119X$.}
\intuicao{A household at zero wealth is not automatically hand-to-mouth at the margin: what matters
is whether the constraint binds \emph{next} period.}
\refbib{Kurlat (2020), \S6.2 (borrowing constraints) and Exercise 6.5; Romer (2012), \S8.6, pp.~393--395.}
\rubrica{
\item Infeasibility: 0.5 pt.
\item KKT conditions with signed multipliers: 1.5 pts. Asserting $c=y$ without the multipliers: max.\ 0.5 pt here.
\item Welfare cost: 1 pt.
\item Both propensities, with the timing argument: 1 pt.
}
\end{sol}

% =============================================================================
\bloco{2}{Labor supply and search}{Kurlat ch.~7; Romer ch.~10 \,\textbullet\, 25 points}

\questao{(10 points) Objective items.}

\textbf{a)} (3 points) Let $u=\ln c-\psi\frac{h^{1+1/\varepsilon}}{1+1/\varepsilon}$ with $c=wh+\pi$, evaluated at $\pi=0$, $\varepsilon=1/2$. The
Marshallian, Hicksian (compensated) and Frisch elasticities of hours with respect to $w$ are:
\begin{alternativas}
\item $(0,\ 1/3,\ 1/2)$
\item $(1/2,\ 1/2,\ 1/2)$
\item $(0,\ 1/2,\ 1/3)$
\item $(-1/3,\ 0,\ 1/2)$
\end{alternativas}
\gabarito{A}
\begin{sol}
\passo{FOC}{$\psi h^{1/\varepsilon}=w/c$. Marshallian: at $\pi=0$, $c=wh$ gives $\psi h^{1+1/\varepsilon}=1$: hours do not depend on $w$.}
\passo{Frisch}{Holding the marginal utility of wealth $1/c$ fixed, $h\propto w^{\varepsilon}$: elasticity $1/2$.}
\passo{Hicksian}{Implicit differentiation of $\psi h^{\phi}(wh+\pi)=w$ ($\phi=1/\varepsilon$) gives
$\partial h/\partial\pi=-1/[w(1+\phi)]$ at $\pi=0$. Slutsky: $\partial h^c/\partial w=\partial h/\partial w-h\,\partial h/\partial\pi$, so the
compensated elasticity is $1/(1+\phi)=\varepsilon/(1+\varepsilon)=1/3$. Frisch $\ge$ Hicksian $\ge$ Marshallian.}
\refbib{Kurlat (2020), \S7.2, eq.\ (7.2.4), pp.~134--135, and footnote 8, p.~149.}
\end{sol}

\medskip
\textbf{b)} (3 points) Prescott's calculation with $u=\ln c+\alpha\ln(1-h)$, $Y=K^{\vartheta}(hN)^{1-\vartheta}$, $\vartheta=1/3$, a
labor tax wedge $\tau$ rebated lump-sum. Calibrate $\alpha$ on the US ($h=0.40$, $\tau=0.40$, $c/y=0.75$)
and predict Brazil's hours ($\tau=0.50$, $c/y=0.80$):
\begin{alternativas}
\item $0.400$
\item $0.357$
\item $0.342$
\item $0.333$
\end{alternativas}
\gabarito{C}
\begin{sol}
\passo{FOC}{$\alpha c/(1-h)=(1-\tau)w$ and $w=(1-\vartheta)y/h$ give $\dfrac h{1-h}=\dfrac{(1-\tau)(1-\vartheta)}{\alpha\,c/y}$.}
\passo{calibration}{US: $\frac{0.4}{0.6}=\frac{0.6\cdot2/3}{0.75\alpha}\Rightarrow\alpha=0.8$. Brazil:
$\frac h{1-h}=\frac{0.5\cdot2/3}{0.8\cdot0.8}=0.521$, $h=0.342$. Option B keeps $c/y=0.75$.}
\refbib{Kurlat (2020), \S7.3 and Exercise 7.5, pp.~137--140 and 148--149.}
\end{sol}

\medskip
\textbf{c)} (2 points, T or F) With matching function $m=\mu U^{1-\alpha}V^{\alpha}$, the search equilibrium with
bargained wages is efficient if and only if the workers' bargaining power equals $\alpha$.
\gabarito{F}
\begin{sol}
\passo{Hosios}{Efficiency requires the workers' power to equal the elasticity of matches with
respect to \emph{unemployment}, $1-\alpha$ (Question 4(c)).}
\refbib{Romer (2012), \S10.7, pp.~497--498, and problem 10.17.}
\end{sol}

\medskip
\textbf{d)} (2 points, T or F) In the steady-state search model with $w=(1-\phi)z+\phi(y+\kappa\theta)$, if
unemployment income $z$ and the vacancy cost $\kappa$ are proportional to productivity $y$, a permanent
rise in $y$ leaves unemployment unchanged.
\gabarito{T}
\begin{sol}
\passo{homogeneity}{The job-creation condition $(1-\phi)(y-z)q(\theta)=\kappa(r+s+\phi f(\theta))$ is homogeneous of
degree one in $(y,z,\kappa)$; with $z=\zeta y$, $\kappa=ky$, $y$ cancels and $\theta$, hence $u$, is unchanged.}
\refbib{Kurlat (2020), \S7.5, p.~146 (Shimer's critique).}
\end{sol}

% -----------------------------------------------------------------------------
\questao{(15 points) Discursive. The search-and-matching model in steady state (monthly). Labor
force $1$; unemployment $u$, vacancies $v$, tightness $\theta=v/u$; matches $m=\mu u^{1-\alpha}v^{\alpha}$;
job-finding rate $f(\theta)=m/u=\mu\theta^{\alpha}$; vacancy-filling rate $q(\theta)=m/v=\mu\theta^{\alpha-1}$;
separations at rate $s$. A filled job produces $y$ and pays $w$; a vacancy costs $\kappa$ per month;
entry is free. Take the wage rule $w=(1-\phi)z+\phi(y+\kappa\theta)$ as given (it comes from Nash
bargaining, $\phi$ = workers' power). Use one fact: a flow $x$ that lasts until a Poisson event with
rate $\varrho$ is worth $x/(r+\varrho)$. Numbers: $r=0.004$, $s=0.03$, $\mu=0.3$, $\alpha=0.5$, $\phi=0.5$, $y=1$, $z=0.7$.}

\textbf{a)} (3 points) Derive the Beveridge curve from flow balance and prove that, in the $(u,v)$
plane, it is downward sloping and convex.
\resolespaco
\begin{sol}
\passo{flows}{$\dot u=s(1-u)-f(\theta)u$. Steady state: $u=\dfrac s{s+f(\theta)}$.}
\passo{in $(u,v)$}{$s(1-u)=\mu u^{1-\alpha}v^{\alpha}\Rightarrow v=C(1-u)^{1/\alpha}u^{-(1-\alpha)/\alpha}$, $C=(s/\mu)^{1/\alpha}$.}
\passo{shape}{Both factors are positive, decreasing and convex in $u$ ($1/\alpha\ge1$). For a product
$fg$ of such functions $(fg)'=f'g+fg'<0$ and $(fg)''=f''g+2f'g'+fg''>0$, since $f'g'>0$.}
\refbib{Kurlat (2020), \S7.1 and \S7.5, eq.\ (7.5.2); Romer (2012), \S10.6.}
\rubrica{
\item Steady-state flow balance: 1 pt.
\item Curve in $(u,v)$: 1 pt.
\item Slope and convexity proved: 1 pt.
}
\end{sol}

\textbf{b)} (4 points) Derive the job-creation condition from free entry, combine it with the wage
rule, and prove that tightness $\theta^*$ exists and is unique. Calibrate $\kappa$ so that $\theta^*=1$ and compute
$u^*$ and $v^*$. Sign $d\theta^*/dy$, $d\theta^*/dz$, $d\theta^*/d\phi$ and $d\theta^*/d\mu$.
\resolespaco
\begin{sol}
\passo{filled job}{A flow $y-w$ until separation: $J=\dfrac{y-w}{r+s}$.}
\passo{vacancy}{A flow $-\kappa$ until filled (rate $q$), then worth $J$: its value is $\dfrac{-\kappa+qJ}{r+q}$. Free entry
sets it to zero: $\kappa/q(\theta)=J$ (expected recruiting cost $=$ value of a filled job).}
\passo{job creation}{With the wage rule, $y-w=(1-\phi)(y-z)-\phi\kappa\theta$, so
\[(1-\phi)(y-z)=\kappa\Big[\frac{r+s}{q(\theta)}+\phi\theta\Big]\iff(1-\phi)(y-z)q(\theta)=\kappa\big(r+s+\phi f(\theta)\big).\]}
\passo{existence and uniqueness}{The bracket is continuous and strictly increasing in $\theta$
($1/q=\theta^{1-\alpha}/\mu$), equals $0$ at $\theta=0$ and tends to $\infty$: if $y>z$ there is exactly one $\theta^*>0$.}
\passo{calibration}{At $\theta=1$: $\kappa=\dfrac{(1-\phi)(y-z)\mu}{r+s+\phi\mu}=\dfrac{0.045}{0.184}=0.2446$;
$f=0.3$, $u^*=0.03/0.33=9.09\%$, $v^*=\theta^*u^*=9.09\%$.}
\passo{comparative statics}{Write $G(\theta)\equiv\kappa[(r+s)\theta^{1-\alpha}/\mu+\phi\theta]-(1-\phi)(y-z)=0$, $G_\theta>0$.
$G_y<0\Rightarrow d\theta^*/dy>0$; $G_z>0\Rightarrow d\theta^*/dz<0$; $G_\phi=\kappa\theta+y-z>0\Rightarrow d\theta^*/d\phi<0$;
$G_\mu<0\Rightarrow d\theta^*/d\mu>0$ (implicit function theorem: $d\theta/dx=-G_x/G_\theta$).}
\refbib{Romer (2012), \S10.6, pp.~486--492 (present values instead of asset equations); Kurlat (2020), \S7.5, eq.\ (7.5.1).}
\rubrica{
\item $J$ and the free-entry condition as present values: 1 pt.
\item Job-creation condition with the wage: 1 pt.
\item Existence and uniqueness: 1 pt. Graph only: 0.5 pt.
\item Calibration and the four signs: 1 pt.
}
\end{sol}

\textbf{c)} (4 points) Efficiency with $r\to0$. A planner chooses $\theta$ to maximize steady-state net output
$y(1-u)+zu-\kappa v$ subject to the Beveridge curve. Derive its condition and prove that the market
$\theta$ coincides with it if and only if $\phi=1-\alpha$. With $\kappa$ from (b), compute efficient $\theta$ and $u$, and the
market values when $\phi=0.7$. Explain the externalities.
\resolespaco
\begin{sol}
\passo{objective}{$W(\theta)=y-(y-z)u(\theta)-\kappa\theta u(\theta)$, $u=\frac s{s+f}$, $u'(\theta)=-\frac{u\,f'(\theta)}{s+f}$, $f'=\alpha q$.}
\passo{planner FOC}{$W'=-(y-z+\kappa\theta)u'-\kappa u=0\Rightarrow(y-z+\kappa\theta)\frac{\alpha q}{s+f}=\kappa\Rightarrow
\alpha(y-z+\kappa\theta)=\kappa\frac{s}{q}+\kappa\theta$, i.e.\
\[\alpha(y-z)=\kappa\frac{s}{q(\theta)}+(1-\alpha)\kappa\theta.\]}
\passo{market at $r=0$}{From (b): $(1-\phi)(y-z)=\kappa\frac{s}{q(\theta)}+\phi\kappa\theta$. Both right-hand sides are
increasing in $\theta$, so the two solutions coincide for all parameters iff $1-\phi=\alpha$ and $\phi=1-\alpha$: Hosios.}
\passo{numbers}{Efficient: $\theta=1.024$, $u=8.99\%$ (the market with $\phi=0.5$ reaches it). With $\phi=0.7$:
$\theta=0.432$, $u=13.21\%$.}
\passo{externalities}{An extra vacancy raises $f$ (helps the unemployed) and lowers $q$ (congests other
firms). Entry is efficient when the firm's share of the surplus, $1-\phi$, equals its contribution to
matching, $\alpha$.}
\refbib{Romer (2012), \S10.7, pp.~497--498, and problem 10.17, p.~512.}
\rubrica{
\item Planner objective and FOC: 1.5 pts.
\item Comparison and the Hosios condition: 1.5 pts. Stating $\phi=1-\alpha$ without the comparison: max.\ 1 pt in the item.
\item Numbers and externalities: 1 pt.
}
\end{sol}

\textbf{d)} (4 points) With $\alpha=1/2$ the job-creation condition is a quadratic in $x=\sqrt\theta$: solve it. Compute
$\theta^*$ and $u^*$ after (i) mismatch lowers $\mu$ to $0.225$ and (ii) more generous unemployment insurance
raises $z$ to $0.8$. Derive the elasticity of $\theta^*$ with respect to $y$ at the calibrated point, and of $u^*$.
What does the formula say about Shimer's puzzle?
\resolespaco
\begin{sol}
\passo{quadratic}{Multiply $(1-\phi)(y-z)\mu x^{-1}=\kappa(r+s+\phi\mu x)$ by $x$:
$\kappa\phi\mu x^2+\kappa(r+s)x-(1-\phi)(y-z)\mu=0$, with positive root
$x=\dfrac{-\kappa(r+s)+\sqrt{\kappa^2(r+s)^2+4\kappa\phi\mu^2(1-\phi)(y-z)}}{2\kappa\phi\mu}$.}
\passo{shocks}{(i) $\mu=0.225$: $\theta^*=0.935$, $u^*=12.12\%$; both the Beveridge curve shifts out and $\theta$ falls.
(ii) $z=0.8$: $\theta^*=0.637$, $u^*=11.14\%$.}
\passo{elasticity}{Implicit differentiation of the job-creation condition at $\theta=1$, with $\kappa$ from (b):
\[\frac{d\ln\theta^*}{d\ln y}=\frac{y}{y-z}\cdot\frac1{(1-\alpha)+\alpha\phi f/(r+s+\phi f)}=3.33\cdot\frac{1}{0.908}=3.67,\]
and $\frac{d\ln u^*}{d\ln\theta}=-\alpha(1-u^*)=-0.455$, so $\frac{d\ln u^*}{d\ln y}=-1.67$.}
\passo{Shimer}{The response is governed by $y/(y-z)$: with a large surplus ($z$ far below $y$) the
model generates little volatility in $\theta$ and $u$; matching the data requires a small surplus ($z$
close to $y$) or rigid wages.}
\refbib{Kurlat (2020), \S7.5, p.~146; Romer (2012), problems 10.12--10.15.}
\rubrica{
\item Quadratic and its root: 1 pt.
\item Both shocks: 1 pt.
\item Elasticity derived by implicit differentiation: 1.5 pts.
\item Shimer reading: 0.5 pt.
}
\end{sol}

\end{document}
